\subsection*{Monotonicity of Sequences}

\begin{exposition}
Monotonicity is an order condition on successive terms of a sequence. It is
independent of convergence by itself, but when combined with boundedness it
becomes one of the fundamental convergence mechanisms in real analysis.
\end{exposition}

% ---------------------------------------------------------
% Definition --- Increasing Sequence
% ---------------------------------------------------------

\begin{definitionbox}{Definition (Increasing Sequence)}
\begin{definition}[Increasing Sequence]
\label{def:increasing-sequence}
Let $(x_n)$ be a real sequence. The sequence $(x_n)$ is
\textbf{increasing} if
\[
  x_n \le x_{n+1}
  \qquad
  \text{for every } n\in\mathbb{N}.
\]
\end{definition}
\end{definitionbox}

\begin{remark*}[Standard quantified statement]
\[
(x_n)\text{ is increasing}
\Longleftrightarrow
\forall n\in\mathbb{N}\;(x_n\le x_{n+1}).
\]
\end{remark*}

\begin{remark*}[Definition predicate reading]
\[
\operatorname{MonotoneIncreasingSequence}(x_n)
\coloneqq
\forall n\in\mathbb{N}\;(x_n\le x_{n+1}).
\]
\end{remark*}

\begin{remark*}[Negated quantified statement]
\[
(x_n)\text{ is not increasing}
\Longleftrightarrow
\exists n\in\mathbb{N}\;(x_{n+1}<x_n).
\]
\end{remark*}

\begin{remark*}[Negation predicate reading]
\[
\neg\operatorname{MonotoneIncreasingSequence}(x_n)
\Longleftrightarrow
\exists n\in\mathbb{N}\;(x_{n+1}<x_n).
\]
\end{remark*}

\begin{remark*}[Interpretation]
An increasing sequence may stay flat. The term ``increasing'' is therefore
used here in the non-strict sense: each term is at least as large as the one
before it.
\end{remark*}

\begin{dependencies}
\begin{itemize}
  \item \hyperref[def:sequence]{Sequence}
\end{itemize}
\end{dependencies}

% ---------------------------------------------------------
% Definition --- Decreasing Sequence
% ---------------------------------------------------------

\begin{definitionbox}{Definition (Decreasing Sequence)}
\begin{definition}[Decreasing Sequence]
\label{def:decreasing-sequence}
Let $(x_n)$ be a real sequence. The sequence $(x_n)$ is
\textbf{decreasing} if
\[
  x_{n+1} \le x_n
  \qquad
  \text{for every } n\in\mathbb{N}.
\]
\end{definition}
\end{definitionbox}

\begin{remark*}[Standard quantified statement]
\[
(x_n)\text{ is decreasing}
\Longleftrightarrow
\forall n\in\mathbb{N}\;(x_{n+1}\le x_n).
\]
\end{remark*}

\begin{remark*}[Definition predicate reading]
\[
\operatorname{MonotoneDecreasingSequence}(x_n)
\coloneqq
\forall n\in\mathbb{N}\;(x_{n+1}\le x_n).
\]
\end{remark*}

\begin{remark*}[Negated quantified statement]
\[
(x_n)\text{ is not decreasing}
\Longleftrightarrow
\exists n\in\mathbb{N}\;(x_n<x_{n+1}).
\]
\end{remark*}

\begin{remark*}[Negation predicate reading]
\[
\neg\operatorname{MonotoneDecreasingSequence}(x_n)
\Longleftrightarrow
\exists n\in\mathbb{N}\;(x_n<x_{n+1}).
\]
\end{remark*}

\begin{remark*}[Interpretation]
A decreasing sequence may stay flat. The order condition says the sequence
never rises as the index advances.
\end{remark*}

\begin{dependencies}
\begin{itemize}
  \item \hyperref[def:sequence]{Sequence}
\end{itemize}
\end{dependencies}

% ---------------------------------------------------------
% Definition --- Strictly Increasing Sequence
% ---------------------------------------------------------

\begin{definitionbox}{Definition (Strictly Increasing Sequence)}
\begin{definition}[Strictly Increasing Sequence]
\label{def:strictly-increasing-sequence}
Let $(x_n)$ be a real sequence. The sequence $(x_n)$ is
\textbf{strictly increasing} if
\[
  x_n < x_{n+1}
  \qquad
  \text{for every } n\in\mathbb{N}.
\]
\end{definition}
\end{definitionbox}

\begin{remark*}[Standard quantified statement]
\[
(x_n)\text{ is strictly increasing}
\Longleftrightarrow
\forall n\in\mathbb{N}\;(x_n<x_{n+1}).
\]
\end{remark*}

\begin{remark*}[Definition predicate reading]
\[
\operatorname{StrictlyIncreasingSequence}(x_n)
\coloneqq
\forall n\in\mathbb{N}\;(x_n<x_{n+1}).
\]
\end{remark*}

\begin{remark*}[Negated quantified statement]
\[
(x_n)\text{ is not strictly increasing}
\Longleftrightarrow
\exists n\in\mathbb{N}\;(x_{n+1}\le x_n).
\]
\end{remark*}

\begin{remark*}[Negation predicate reading]
\[
\neg\operatorname{StrictlyIncreasingSequence}(x_n)
\Longleftrightarrow
\exists n\in\mathbb{N}\;(x_{n+1}\le x_n).
\]
\end{remark*}

\begin{remark*}[Interpretation]
Strict increase excludes equality between consecutive terms. Every strictly
increasing sequence is increasing, but the converse need not hold.
\end{remark*}

\begin{dependencies}
\begin{itemize}
  \item \hyperref[def:sequence]{Sequence}
  \item \hyperref[def:increasing-sequence]{Increasing Sequence}
\end{itemize}
\end{dependencies}

% ---------------------------------------------------------
% Definition --- Strictly Decreasing Sequence
% ---------------------------------------------------------

\begin{definitionbox}{Definition (Strictly Decreasing Sequence)}
\begin{definition}[Strictly Decreasing Sequence]
\label{def:strictly-decreasing-sequence}
Let $(x_n)$ be a real sequence. The sequence $(x_n)$ is
\textbf{strictly decreasing} if
\[
  x_{n+1} < x_n
  \qquad
  \text{for every } n\in\mathbb{N}.
\]
\end{definition}
\end{definitionbox}

\begin{remark*}[Standard quantified statement]
\[
(x_n)\text{ is strictly decreasing}
\Longleftrightarrow
\forall n\in\mathbb{N}\;(x_{n+1}<x_n).
\]
\end{remark*}

\begin{remark*}[Definition predicate reading]
\[
\operatorname{StrictlyDecreasingSequence}(x_n)
\coloneqq
\forall n\in\mathbb{N}\;(x_{n+1}<x_n).
\]
\end{remark*}

\begin{remark*}[Negated quantified statement]
\[
(x_n)\text{ is not strictly decreasing}
\Longleftrightarrow
\exists n\in\mathbb{N}\;(x_n\le x_{n+1}).
\]
\end{remark*}

\begin{remark*}[Negation predicate reading]
\[
\neg\operatorname{StrictlyDecreasingSequence}(x_n)
\Longleftrightarrow
\exists n\in\mathbb{N}\;(x_n\le x_{n+1}).
\]
\end{remark*}

\begin{remark*}[Interpretation]
Strict decrease excludes equality between consecutive terms. Every strictly
decreasing sequence is decreasing, but the converse need not hold.
\end{remark*}

\begin{dependencies}
\begin{itemize}
  \item \hyperref[def:sequence]{Sequence}
  \item \hyperref[def:decreasing-sequence]{Decreasing Sequence}
\end{itemize}
\end{dependencies}

% ---------------------------------------------------------
% Definition --- Monotone Sequence
% ---------------------------------------------------------

\begin{definitionbox}{Definition (Monotone Sequence)}
\begin{definition}[Monotone Sequence]
\label{def:monotone-sequence}
Let $(x_n)$ be a real sequence. The sequence $(x_n)$ is
\textbf{monotone} if it is increasing or decreasing.
\end{definition}
\end{definitionbox}

\begin{remark*}[Standard quantified statement]
\[
\begin{aligned}
&(x_n)\text{ is monotone}\\
&\quad\Longleftrightarrow\quad
\bigl[\forall n\in\mathbb{N}\;(x_n\le x_{n+1})\bigr]
\lor
\bigl[\forall n\in\mathbb{N}\;(x_{n+1}\le x_n)\bigr].
\end{aligned}
\]
\end{remark*}

\begin{remark*}[Definition predicate reading]
\[
\operatorname{MonotoneSequence}(x_n)
\coloneqq
\operatorname{MonotoneIncreasingSequence}(x_n)
\lor
\operatorname{MonotoneDecreasingSequence}(x_n).
\]
\end{remark*}

\begin{remark*}[Negated quantified statement]
\[
\begin{aligned}
&(x_n)\text{ is not monotone}\\
&\quad\Longleftrightarrow\quad
\bigl[\exists n\in\mathbb{N}\;(x_{n+1}<x_n)\bigr]
\land
\bigl[\exists m\in\mathbb{N}\;(x_m<x_{m+1})\bigr].
\end{aligned}
\]
\end{remark*}

\begin{remark*}[Negation predicate reading]
\[
\neg\operatorname{MonotoneSequence}(x_n)
\Longleftrightarrow
\neg\operatorname{MonotoneIncreasingSequence}(x_n)
\land
\neg\operatorname{MonotoneDecreasingSequence}(x_n).
\]
\end{remark*}

\begin{remark*}[Interpretation]
Monotonicity means the sequence has one global direction. It may move upward,
move downward, or stay constant, but it cannot reverse direction.
\end{remark*}

\begin{dependencies}
\begin{itemize}
  \item \hyperref[def:increasing-sequence]{Increasing Sequence}
  \item \hyperref[def:decreasing-sequence]{Decreasing Sequence}
\end{itemize}
\end{dependencies}

% ---------------------------------------------------------
% Definition --- Eventually Increasing Sequence
% ---------------------------------------------------------

\begin{definitionbox}{Definition (Eventually Increasing Sequence)}
\begin{definition}[Eventually Increasing Sequence]
\label{def:eventually-increasing-sequence}
Let $(x_n)$ be a real sequence. The sequence $(x_n)$ is
\textbf{eventually increasing} if there exists $N\in\mathbb{N}$ such that
\[
  x_n \le x_{n+1}
  \qquad
  \text{for every } n\ge N.
\]
\end{definition}
\end{definitionbox}

\begin{remark*}[Standard quantified statement]
\[
(x_n)\text{ is eventually increasing}
\Longleftrightarrow
\exists N\in\mathbb{N}\;\forall n\ge N\;(x_n\le x_{n+1}).
\]
\end{remark*}

\begin{remark*}[Definition predicate reading]
\[
\operatorname{EventuallyIncreasingSequence}(x_n)
\coloneqq
\exists N\in\mathbb{N}\;\forall n\ge N\;(x_n\le x_{n+1}).
\]
\end{remark*}

\begin{remark*}[Negated quantified statement]
\[
(x_n)\text{ is not eventually increasing}
\Longleftrightarrow
\forall N\in\mathbb{N}\;\exists n\ge N\;(x_{n+1}<x_n).
\]
\end{remark*}

\begin{remark*}[Negation predicate reading]
\[
\neg\operatorname{EventuallyIncreasingSequence}(x_n)
\Longleftrightarrow
\forall N\in\mathbb{N}\;\exists n\ge N\;(x_{n+1}<x_n).
\]
\end{remark*}

\begin{remark*}[Interpretation]
Eventual increase allows finitely many early reversals. Only the tail of the
sequence must move in a nondecreasing direction.
\end{remark*}

\begin{dependencies}
\begin{itemize}
  \item \hyperref[def:m-tail]{$M$-Tail of a Sequence}
  \item \hyperref[def:increasing-sequence]{Increasing Sequence}
\end{itemize}
\end{dependencies}

% ---------------------------------------------------------
% Definition --- Eventually Decreasing Sequence
% ---------------------------------------------------------

\begin{definitionbox}{Definition (Eventually Decreasing Sequence)}
\begin{definition}[Eventually Decreasing Sequence]
\label{def:eventually-decreasing-sequence}
Let $(x_n)$ be a real sequence. The sequence $(x_n)$ is
\textbf{eventually decreasing} if there exists $N\in\mathbb{N}$ such that
\[
  x_{n+1} \le x_n
  \qquad
  \text{for every } n\ge N.
\]
\end{definition}
\end{definitionbox}

\begin{remark*}[Standard quantified statement]
\[
(x_n)\text{ is eventually decreasing}
\Longleftrightarrow
\exists N\in\mathbb{N}\;\forall n\ge N\;(x_{n+1}\le x_n).
\]
\end{remark*}

\begin{remark*}[Definition predicate reading]
\[
\operatorname{EventuallyDecreasingSequence}(x_n)
\coloneqq
\exists N\in\mathbb{N}\;\forall n\ge N\;(x_{n+1}\le x_n).
\]
\end{remark*}

\begin{remark*}[Negated quantified statement]
\[
(x_n)\text{ is not eventually decreasing}
\Longleftrightarrow
\forall N\in\mathbb{N}\;\exists n\ge N\;(x_n<x_{n+1}).
\]
\end{remark*}

\begin{remark*}[Negation predicate reading]
\[
\neg\operatorname{EventuallyDecreasingSequence}(x_n)
\Longleftrightarrow
\forall N\in\mathbb{N}\;\exists n\ge N\;(x_n<x_{n+1}).
\]
\end{remark*}

\begin{remark*}[Interpretation]
Eventual decrease allows the initial terms to behave irregularly. The order
condition begins only after some finite threshold.
\end{remark*}

\begin{dependencies}
\begin{itemize}
  \item \hyperref[def:m-tail]{$M$-Tail of a Sequence}
  \item \hyperref[def:decreasing-sequence]{Decreasing Sequence}
\end{itemize}
\end{dependencies}

% ---------------------------------------------------------
% Definition --- Eventually Monotone Sequence
% ---------------------------------------------------------

\begin{definitionbox}{Definition (Eventually Monotone Sequence)}
\begin{definition}[Eventually Monotone Sequence]
\label{def:eventually-monotone-sequence}
Let $(x_n)$ be a real sequence. The sequence $(x_n)$ is
\textbf{eventually monotone} if it is eventually increasing or eventually
decreasing.
\end{definition}
\end{definitionbox}

\begin{remark*}[Standard quantified statement]
\[
\begin{aligned}
&(x_n)\text{ is eventually monotone}\\
&\quad\Longleftrightarrow\quad
\bigl[\exists N\in\mathbb{N}\;\forall n\ge N\;(x_n\le x_{n+1})\bigr]
\lor
\bigl[\exists N\in\mathbb{N}\;\forall n\ge N\;(x_{n+1}\le x_n)\bigr].
\end{aligned}
\]
\end{remark*}

\begin{remark*}[Definition predicate reading]
\[
\operatorname{EventuallyMonotoneSequence}(x_n)
\coloneqq
\operatorname{EventuallyIncreasingSequence}(x_n)
\lor
\operatorname{EventuallyDecreasingSequence}(x_n).
\]
\end{remark*}

\begin{remark*}[Negated quantified statement]
\[
\begin{aligned}
&(x_n)\text{ is not eventually monotone}\\
&\quad\Longleftrightarrow\quad
\bigl[\forall N\in\mathbb{N}\;\exists n\ge N\;(x_{n+1}<x_n)\bigr]
\land
\bigl[\forall N\in\mathbb{N}\;\exists n\ge N\;(x_n<x_{n+1})\bigr].
\end{aligned}
\]
\end{remark*}

\begin{remark*}[Negation predicate reading]
\[
\neg\operatorname{EventuallyMonotoneSequence}(x_n)
\Longleftrightarrow
\neg\operatorname{EventuallyIncreasingSequence}(x_n)
\land
\neg\operatorname{EventuallyDecreasingSequence}(x_n).
\]
\end{remark*}

\begin{remark*}[Interpretation]
Eventual monotonicity is the tail version of monotonicity. It is often the
right hypothesis in applications, because finitely many early terms do not
affect convergence.
\end{remark*}

\begin{dependencies}
\begin{itemize}
  \item \hyperref[def:eventually-increasing-sequence]{Eventually Increasing Sequence}
  \item \hyperref[def:eventually-decreasing-sequence]{Eventually Decreasing Sequence}
  \item \hyperref[def:m-tail]{Convergence of a Tail}
\end{itemize}
\end{dependencies}

% ---------------------------------------------------------
% Theorem --- Monotone Convergence Theorem
% ---------------------------------------------------------

\begin{theorembox}{Theorem (Monotone Convergence Theorem)}
\begin{theorem}[Monotone Convergence Theorem]
\label{thm:monotone-convergence-theorem}
Let $(x_n)$ be a real sequence.
\begin{enumerate}[label=(\alph*)]
  \item If $(x_n)$ is increasing and bounded above, then $(x_n)$ converges and
  \[
    \lim_{n\to\infty}x_n=\sup\{x_n:n\in\mathbb{N}\}.
  \]
  \item If $(x_n)$ is decreasing and bounded below, then $(x_n)$ converges and
  \[
    \lim_{n\to\infty}x_n=\inf\{x_n:n\in\mathbb{N}\}.
  \]
\end{enumerate}

\smallskip
\noindent
\hyperref[prf:monotone-convergence-theorem]{\textit{Go to proof.}}
\end{theorem}
\end{theorembox}

\begin{remark*}[Standard quantified statement]
\[
\begin{aligned}
&\forall (x_n)\subseteq\mathbb{R}\\
&\quad
\Bigl(
  \bigl[\forall n\in\mathbb{N}\;(x_n\le x_{n+1})\bigr]
  \land
  \bigl[\exists M\in\mathbb{R}\;\forall n\in\mathbb{N}\;(x_n\le M)\bigr]
\Bigr)\\
&\quad\Longrightarrow\quad
\exists L\in\mathbb{R}\;\Bigl(
  \lim_{n\to\infty}x_n=L
  \land
  L=\sup\{x_n:n\in\mathbb{N}\}
\Bigr),\\
&\quad
\Bigl(
  \bigl[\forall n\in\mathbb{N}\;(x_{n+1}\le x_n)\bigr]
  \land
  \bigl[\exists m\in\mathbb{R}\;\forall n\in\mathbb{N}\;(m\le x_n)\bigr]
\Bigr)\\
&\quad\Longrightarrow\quad
\exists L\in\mathbb{R}\;\Bigl(
  \lim_{n\to\infty}x_n=L
  \land
  L=\inf\{x_n:n\in\mathbb{N}\}
\Bigr).
\end{aligned}
\]
\end{remark*}

\begin{remark*}[Theorem predicate reading]
\[
\begin{aligned}
&\operatorname{MonotoneIncreasingSequence}(x_n)
\land
\operatorname{BoundedAboveSequence}(x_n)
\Longrightarrow
\operatorname{ConvergentSequence}\bigl(x_n,\sup\{x_n:n\in\mathbb{N}\}\bigr),\\
&\operatorname{MonotoneDecreasingSequence}(x_n)
\land
\operatorname{BoundedBelowSequence}(x_n)
\Longrightarrow
\operatorname{ConvergentSequence}\bigl(x_n,\inf\{x_n:n\in\mathbb{N}\}\bigr).
\end{aligned}
\]
\end{remark*}

\begin{remark*}[Negated quantified statement]
\[
\begin{aligned}
&\exists (x_n)\subseteq\mathbb{R}\;\Bigl(
  \bigl[\forall n\in\mathbb{N}\;(x_n\le x_{n+1})\bigr]
  \land
  \bigl[\exists M\in\mathbb{R}\;\forall n\in\mathbb{N}\;(x_n\le M)\bigr]\\
&\qquad\qquad\qquad\qquad
  \land
  \text{$(x_n)$ does not converge}
\Bigr),\\
&\exists (x_n)\subseteq\mathbb{R}\;\Bigl(
  \bigl[\forall n\in\mathbb{N}\;(x_{n+1}\le x_n)\bigr]
  \land
  \bigl[\exists m\in\mathbb{R}\;\forall n\in\mathbb{N}\;(m\le x_n)\bigr]\\
&\qquad\qquad\qquad\qquad
  \land
  \text{$(x_n)$ does not converge}
\Bigr).
\end{aligned}
\]
\end{remark*}

\begin{remark*}[Negation predicate reading]
\[
\begin{aligned}
&\operatorname{MonotoneIncreasingSequence}(x_n)
\land
\operatorname{BoundedAboveSequence}(x_n)
\land
\neg\exists L\in\mathbb{R}\;\operatorname{ConvergentSequence}(x_n,L),\\
&\operatorname{MonotoneDecreasingSequence}(x_n)
\land
\operatorname{BoundedBelowSequence}(x_n)
\land
\neg\exists L\in\mathbb{R}\;\operatorname{ConvergentSequence}(x_n,L).
\end{aligned}
\]
\end{remark*}

\begin{remark*}[Failure modes]
The theorem would fail only if a one-directional bounded sequence could avoid
settling to a real limit. Completeness of $\mathbb{R}$ rules this out: an
increasing bounded-above sequence is forced toward the least upper bound of
its range, while a decreasing bounded-below sequence is forced toward the
greatest lower bound of its range.
\end{remark*}

\begin{remark*}[Failure mode decomposition]
\[
\underbrace{\forall n\in\mathbb{N}\;(x_n\le x_{n+1})}_{\text{one direction}}
\quad
\underbrace{\exists M\in\mathbb{R}\;\forall n\in\mathbb{N}\;(x_n\le M)}_{\text{ceiling}}
\quad
\underbrace{\text{$(x_n)$ does not converge}}_{\text{forbidden by completeness}}.
\]
The decreasing case is dual: reverse the inequalities and replace the
ceiling by a floor.
\end{remark*}

\begin{remark*}[Interpretation]
Boundedness prevents escape, and monotonicity prevents oscillation. Together
they force convergence. The theorem is one of the standard forms in which the
completeness of the real numbers becomes a sequence theorem.
\end{remark*}

\begin{dependencies}
\begin{itemize}
  \item \hyperref[def:convergent-sequence]{Convergence of a Sequence}
  \item \hyperref[def:increasing-sequence]{Increasing Sequence}
  \item \hyperref[def:decreasing-sequence]{Decreasing Sequence}
  \item \hyperref[def:sequence-bounded-above]{Sequence Bounded Above}
  \item \hyperref[def:sequence-bounded-below]{Sequence Bounded Below}
  \item \hyperref[def:supremum]{Supremum}
  \item \hyperref[def:infimum]{Infimum}
\end{itemize}
\end{dependencies}

% ---------------------------------------------------------
% Theorem --- Strict Monotonicity Implies Monotonicity
% ---------------------------------------------------------

\begin{theorembox}{Theorem (Strict Monotonicity Implies Monotonicity)}
\begin{theorem}[Strict Monotonicity Implies Monotonicity]
\label{thm:strict-monotonicity-implies-monotonicity}
Let $(x_n)$ be a real sequence.
\begin{enumerate}[label=(\alph*)]
  \item If $(x_n)$ is strictly increasing, then $(x_n)$ is increasing.
  \item If $(x_n)$ is strictly decreasing, then $(x_n)$ is decreasing.
\end{enumerate}

\smallskip
\noindent
\hyperref[prf:strict-monotonicity-implies-monotonicity]{\textit{Go to proof.}}
\end{theorem}
\end{theorembox}

\begin{remark*}[Standard quantified statement]
\[
\begin{aligned}
&\forall (x_n)\subseteq\mathbb{R}\\
&\quad\Bigl(
  \bigl[\forall n\in\mathbb{N}\;(x_n<x_{n+1})\bigr]
  \Longrightarrow
  \bigl[\forall n\in\mathbb{N}\;(x_n\le x_{n+1})\bigr]
\Bigr),\\
&\quad\Bigl(
  \bigl[\forall n\in\mathbb{N}\;(x_{n+1}<x_n)\bigr]
  \Longrightarrow
  \bigl[\forall n\in\mathbb{N}\;(x_{n+1}\le x_n)\bigr]
\Bigr).
\end{aligned}
\]
\end{remark*}

\begin{remark*}[Theorem predicate reading]
\[
\begin{aligned}
&\operatorname{StrictlyIncreasingSequence}(x_n)
\Longrightarrow
\operatorname{MonotoneIncreasingSequence}(x_n),\\
&\operatorname{StrictlyDecreasingSequence}(x_n)
\Longrightarrow
\operatorname{MonotoneDecreasingSequence}(x_n).
\end{aligned}
\]
\end{remark*}

\begin{remark*}[Interpretation]
Strict monotonicity is stronger than non-strict monotonicity because every
strict inequality is also a weak inequality.
\end{remark*}

\begin{dependencies}
\begin{itemize}
  \item \hyperref[def:strictly-increasing-sequence]{Strictly Increasing Sequence}
  \item \hyperref[def:strictly-decreasing-sequence]{Strictly Decreasing Sequence}
  \item \hyperref[def:increasing-sequence]{Increasing Sequence}
  \item \hyperref[def:decreasing-sequence]{Decreasing Sequence}
\end{itemize}
\end{dependencies}

% ---------------------------------------------------------
% Theorem --- Bounded Monotone Sequence Equivalences
% ---------------------------------------------------------

\begin{theorembox}{Theorem (Bounded Monotone Sequence Equivalences)}
\begin{theorem}[Bounded Monotone Sequence Equivalences]
\label{thm:bounded-monotone-sequence-equivalences}
Let $(x_n)$ be a real sequence.
\begin{enumerate}[label=(\alph*)]
  \item If $(x_n)$ is increasing, then $(x_n)$ converges if and only if
  $(x_n)$ is bounded above.
  \item If $(x_n)$ is decreasing, then $(x_n)$ converges if and only if
  $(x_n)$ is bounded below.
  \item If $(x_n)$ is monotone, then $(x_n)$ converges if and only if
  $(x_n)$ is bounded.
\end{enumerate}

\smallskip
\noindent
\hyperref[prf:bounded-monotone-sequence-equivalences]{\textit{Go to proof.}}
\end{theorem}
\end{theorembox}

\begin{remark*}[Standard quantified statement]
\[
\begin{aligned}
&\forall (x_n)\subseteq\mathbb{R}\\
&\quad
\bigl[\forall n\in\mathbb{N}\;(x_n\le x_{n+1})\bigr]
\Longrightarrow
\Bigl(
  \exists L\in\mathbb{R}\;(x_n\to L)
  \Longleftrightarrow
  \exists M\in\mathbb{R}\;\forall n\in\mathbb{N}\;(x_n\le M)
\Bigr),\\
&\quad
\bigl[\forall n\in\mathbb{N}\;(x_{n+1}\le x_n)\bigr]
\Longrightarrow
\Bigl(
  \exists L\in\mathbb{R}\;(x_n\to L)
  \Longleftrightarrow
  \exists m\in\mathbb{R}\;\forall n\in\mathbb{N}\;(m\le x_n)
\Bigr).
\end{aligned}
\]
\end{remark*}

\begin{remark*}[Theorem predicate reading]
\[
\begin{aligned}
&\operatorname{MonotoneIncreasingSequence}(x_n)
\Longrightarrow
\Bigl(
\exists L\in\mathbb{R}\;\operatorname{ConvergentSequence}(x_n,L)
\Longleftrightarrow
\operatorname{BoundedAboveSequence}(x_n)
\Bigr),\\
&\operatorname{MonotoneDecreasingSequence}(x_n)
\Longrightarrow
\Bigl(
\exists L\in\mathbb{R}\;\operatorname{ConvergentSequence}(x_n,L)
\Longleftrightarrow
\operatorname{BoundedBelowSequence}(x_n)
\Bigr).
\end{aligned}
\]
\end{remark*}

\begin{remark*}[Interpretation]
For monotone sequences, the only obstruction to real convergence is escape in
the monotone direction. Boundedness removes that obstruction.
\end{remark*}

\begin{dependencies}
\begin{itemize}
  \item \hyperref[thm:monotone-convergence-theorem]{Monotone Convergence Theorem}
  \item \hyperref[def:bounded-sequence]{Bounded Sequence}
  \item \hyperref[def:convergent-sequence]{Convergence of a Sequence}
\end{itemize}
\end{dependencies}

% ---------------------------------------------------------
% Theorem --- Eventually Monotone Convergence Theorem
% ---------------------------------------------------------

\begin{theorembox}{Theorem (Eventually Monotone Convergence Theorem)}
\begin{theorem}[Eventually Monotone Convergence Theorem]
\label{thm:eventually-monotone-convergence-theorem}
Let $(x_n)$ be a real sequence.
\begin{enumerate}[label=(\alph*)]
  \item If $(x_n)$ is eventually increasing and bounded above, then
  $(x_n)$ converges.
  \item If $(x_n)$ is eventually decreasing and bounded below, then
  $(x_n)$ converges.
  \item If $(x_n)$ is eventually monotone and bounded, then $(x_n)$ converges.
\end{enumerate}

\smallskip
\noindent
\hyperref[prf:eventually-monotone-convergence-theorem]{\textit{Go to proof.}}
\end{theorem}
\end{theorembox}

\begin{remark*}[Standard quantified statement]
\[
\begin{aligned}
&\forall (x_n)\subseteq\mathbb{R}\\
&\quad\Bigl(
  \exists N\in\mathbb{N}\;\forall n\ge N\;(x_n\le x_{n+1})
  \land
  \exists M\in\mathbb{R}\;\forall n\in\mathbb{N}\;(x_n\le M)
\Bigr)\\
&\quad\Longrightarrow\quad
\exists L\in\mathbb{R}\;(x_n\to L),\\
&\quad\Bigl(
  \exists N\in\mathbb{N}\;\forall n\ge N\;(x_{n+1}\le x_n)
  \land
  \exists m\in\mathbb{R}\;\forall n\in\mathbb{N}\;(m\le x_n)
\Bigr)\\
&\quad\Longrightarrow\quad
\exists L\in\mathbb{R}\;(x_n\to L).
\end{aligned}
\]
\end{remark*}

\begin{remark*}[Theorem predicate reading]
\[
\begin{aligned}
&\operatorname{EventuallyIncreasingSequence}(x_n)
\land
\operatorname{BoundedAboveSequence}(x_n)
\Longrightarrow
\exists L\in\mathbb{R}\;\operatorname{ConvergentSequence}(x_n,L),\\
&\operatorname{EventuallyDecreasingSequence}(x_n)
\land
\operatorname{BoundedBelowSequence}(x_n)
\Longrightarrow
\exists L\in\mathbb{R}\;\operatorname{ConvergentSequence}(x_n,L).
\end{aligned}
\]
\end{remark*}

\begin{remark*}[Interpretation]
This is the tail version of the Monotone Convergence Theorem. Once the tail is
monotone and bounded in the right direction, finite initial behavior is
irrelevant.
\end{remark*}

\begin{dependencies}
\begin{itemize}
  \item \hyperref[def:eventually-increasing-sequence]{Eventually Increasing Sequence}
  \item \hyperref[def:eventually-decreasing-sequence]{Eventually Decreasing Sequence}
  \item \hyperref[thm:convergence-of-tail]{Convergence of a Tail}
  \item \hyperref[thm:monotone-convergence-theorem]{Monotone Convergence Theorem}
\end{itemize}
\end{dependencies}

% ---------------------------------------------------------
% Theorem --- Unbounded Monotone Divergence
% ---------------------------------------------------------

\begin{theorembox}{Theorem (Unbounded Monotone Divergence)}
\begin{theorem}[Unbounded Monotone Divergence]
\label{thm:unbounded-monotone-divergence}
Let $(x_n)$ be a real sequence.
\begin{enumerate}[label=(\alph*)]
  \item If $(x_n)$ is increasing and is not bounded above, then
  $x_n\to+\infty$.
  \item If $(x_n)$ is decreasing and is not bounded below, then
  $x_n\to-\infty$.
\end{enumerate}

\smallskip
\noindent
\hyperref[prf:unbounded-monotone-divergence]{\textit{Go to proof.}}
\end{theorem}
\end{theorembox}

\begin{remark*}[Standard quantified statement]
\[
\begin{aligned}
&\forall (x_n)\subseteq\mathbb{R}\\
&\quad
\Bigl(
  \bigl[\forall n\in\mathbb{N}\;(x_n\le x_{n+1})\bigr]
  \land
  \bigl[\forall M\in\mathbb{R}\;\exists n\in\mathbb{N}\;(M<x_n)\bigr]
\Bigr)\\
&\quad\Longrightarrow\quad
\forall M\in\mathbb{R}\;\exists N\in\mathbb{N}\;\forall n\ge N\;(M<x_n),\\
&\quad
\Bigl(
  \bigl[\forall n\in\mathbb{N}\;(x_{n+1}\le x_n)\bigr]
  \land
  \bigl[\forall M\in\mathbb{R}\;\exists n\in\mathbb{N}\;(x_n<M)\bigr]
\Bigr)\\
&\quad\Longrightarrow\quad
\forall M\in\mathbb{R}\;\exists N\in\mathbb{N}\;\forall n\ge N\;(x_n<M).
\end{aligned}
\]
\end{remark*}

\begin{remark*}[Theorem predicate reading]
\[
\begin{aligned}
&\operatorname{MonotoneIncreasingSequence}(x_n)
\land
\neg\operatorname{BoundedAboveSequence}(x_n)
\Longrightarrow
\operatorname{DivergesToPositiveInfinity}(x_n),\\
&\operatorname{MonotoneDecreasingSequence}(x_n)
\land
\neg\operatorname{BoundedBelowSequence}(x_n)
\Longrightarrow
\operatorname{DivergesToNegativeInfinity}(x_n).
\end{aligned}
\]
\end{remark*}

\begin{remark*}[Interpretation]
For a monotone sequence, unboundedness has a direction. An increasing sequence
that has no upper bound eventually exceeds every real number; the decreasing
case is the order-dual statement.
\end{remark*}

\begin{dependencies}
\begin{itemize}
  \item \hyperref[def:diverges-to-positive-infinity]{Divergence to Positive Infinity}
  \item \hyperref[def:diverges-to-negative-infinity]{Divergence to Negative Infinity}
  \item \hyperref[def:increasing-sequence]{Increasing Sequence}
  \item \hyperref[def:decreasing-sequence]{Decreasing Sequence}
  \item \hyperref[def:sequence-bounded-above]{Sequence Bounded Above}
  \item \hyperref[def:sequence-bounded-below]{Sequence Bounded Below}
\end{itemize}
\end{dependencies}

% ---------------------------------------------------------
% Theorem --- Algebraic Transformations Preserve Monotonicity
% ---------------------------------------------------------

\begin{theorembox}{Theorem (Algebraic Transformations Preserve Monotonicity)}
\begin{theorem}[Algebraic Transformations Preserve Monotonicity]
\label{thm:algebraic-transformations-preserve-monotonicity}
Let $(x_n)$ be a real sequence and let $c,\alpha\in\mathbb{R}$.
\begin{enumerate}[label=(\alph*)]
  \item If $(x_n)$ is increasing, then $(x_n+c)$ is increasing.
  \item If $(x_n)$ is decreasing, then $(x_n+c)$ is decreasing.
  \item If $(x_n)$ is increasing and $\alpha>0$, then $(\alpha x_n)$ is increasing.
  \item If $(x_n)$ is increasing and $\alpha<0$, then $(\alpha x_n)$ is decreasing.
  \item If $(x_n)$ is increasing, then $(-x_n)$ is decreasing.
  \item If $(x_n)$ is decreasing, then $(-x_n)$ is increasing.
\end{enumerate}

\smallskip
\noindent
\hyperref[prf:algebraic-transformations-preserve-monotonicity]{\textit{Go to proof.}}
\end{theorem}
\end{theorembox}

\begin{remark*}[Standard quantified statement]
\[
\begin{aligned}
&\forall (x_n)\subseteq\mathbb{R}\;\forall c,\alpha\in\mathbb{R}\\
&\quad
\bigl[\forall n\in\mathbb{N}\;(x_n\le x_{n+1})\bigr]
\Longrightarrow
\bigl[\forall n\in\mathbb{N}\;(x_n+c\le x_{n+1}+c)\bigr],\\
&\quad
\bigl(\alpha>0\land\forall n\in\mathbb{N}\;(x_n\le x_{n+1})\bigr)
\Longrightarrow
\bigl[\forall n\in\mathbb{N}\;(\alpha x_n\le \alpha x_{n+1})\bigr],\\
&\quad
\bigl(\alpha<0\land\forall n\in\mathbb{N}\;(x_n\le x_{n+1})\bigr)
\Longrightarrow
\bigl[\forall n\in\mathbb{N}\;(\alpha x_{n+1}\le \alpha x_n)\bigr].
\end{aligned}
\]
\end{remark*}

\begin{remark*}[Theorem predicate reading]
\[
\begin{aligned}
&\operatorname{MonotoneIncreasingSequence}(x_n)
\Longrightarrow
\operatorname{MonotoneIncreasingSequence}(x_n+c),\\
&\operatorname{MonotoneIncreasingSequence}(x_n)\land\alpha>0
\Longrightarrow
\operatorname{MonotoneIncreasingSequence}(\alpha x_n),\\
&\operatorname{MonotoneIncreasingSequence}(x_n)\land\alpha<0
\Longrightarrow
\operatorname{MonotoneDecreasingSequence}(\alpha x_n).
\end{aligned}
\]
\end{remark*}

\begin{remark*}[Interpretation]
Translation preserves order, multiplication by a positive scalar preserves
order, and multiplication by a negative scalar reverses order. Negation is the
special case $\alpha=-1$.
\end{remark*}

\begin{dependencies}
\begin{itemize}
  \item \hyperref[def:pointwise-sum-sequence]{Pointwise Sum of Sequences}
  \item \hyperref[def:scalar-multiple-sequence]{Scalar Multiple of a Sequence}
  \item \hyperref[def:pointwise-negation-sequence]{Pointwise Negation of a Sequence}
  \item \hyperref[def:increasing-sequence]{Increasing Sequence}
  \item \hyperref[def:decreasing-sequence]{Decreasing Sequence}
\end{itemize}
\end{dependencies}
